% Session 3 -- why an "insulator" conducts: chiral edge states in real space.
% Placement: sec. "The quantum Hall state", after the chirality/backscattering
% paragraph; referenced again from the robustness paragraph of sec. 3.7.
%
% LAYOUT NOTE: pure tikzpicture, bare canvas coordinates throughout.
%
% SIGN CONVENTION, checked against the semiclassical drift: electrons (charge
% -e) in a field B out of the page orbit counter-clockwise in the bulk, and the
% confining potential at the boundary gives an E x B drift of +x along the top
% edge and -x along the bottom edge. The edge circulation is therefore
% clockwise, i.e. opposite in sense to the bulk cyclotron motion. Reversing B
% reverses every arrow in the figure at once.
\begin{aqmfigure}{fig:s03:chiraledge}{Chiral edge states, and why they cannot be backscattered.
(a) A quantum Hall sample. In the bulk the electrons move on closed cyclotron orbits and do not go
anywhere, and the system is an insulator with a gap between Landau levels. At the boundary these
orbits cannot close. They skip along the edge, and for a given direction of $B$ they skip in one
direction only, which is opposite on the two edges. Propagation in a single direction is what the
word \textbf{chiral} means, and it carries the current that $\sigma_{xy}=C\,e^2/h$ predicts, even
though the bulk is gapped. (b) The consequence. An impurity at the edge can deflect the mode, but
there is no state that propagates in the opposite direction at the same edge and the same energy for
it to scatter \emph{into}, so the mode goes around the obstacle and continues. Without a
backscattering channel there is no resistance, which is why the Hall plateaus are quantized so
sharply. If we reverse the sign of $B$, all the arrows are reversed at the same time.}

\begin{tikzpicture}[x=1cm, y=1cm]

% =====================================================================
% (a) the sample
% =====================================================================
\def\W{7.2}
\def\Ht{3.9}

\node[aqmpanel, anchor=south west] at (0,\Ht+1.02) {(a) a gapped bulk with conducting edges};

% sample body
\fill[aqmfill] (0,0) rectangle (\W,\Ht);
\draw[aqm, draw=black] (0,0) rectangle (\W,\Ht);

% --- bulk: closed cyclotron orbits (counter-clockwise) ---------------
\foreach \a/\b in {1.00/2.95, 2.60/3.05, 4.30/2.95, 5.90/3.05,
                   6.60/1.95, 3.60/0.95, 5.30/0.95, 6.60/0.95, 1.00/1.95} {
  \draw[aqm, draw=aqmaccent, line width=0.8pt] (\a,\b) circle (0.30);
  \draw[aqm, draw=aqmaccent, -{Stealth[length=4.5pt,width=3.6pt]}, line width=0.8pt]
    ($ (\a,\b) + (-18:0.30) $) arc (-18:22:0.30);
}

% magnetic field symbol
\node[circle, draw=black, fill=white, line width=\aqmthinwidth,
      minimum size=11pt, inner sep=0pt] at (0.62,0.95) {};
\fill[black] (0.62,0.95) circle (1.1pt);
\node[aqmlabsmall, anchor=west] at (0.88,0.95) {$B$ out of page};

\node[aqmlabsmall, align=center, fill=aqmfill, inner sep=2pt]
  at (3.80,1.95)
  {\textbf{bulk:} closed cyclotron orbits,\\[-1pt] gapped, no net transport};

% --- top edge: skipping orbits drifting to the right -----------------
\foreach \i in {0,1,...,9} {
  \draw[aqm, draw=aqmbroken, line width=1.0pt]
    ({0.40+\i*0.66},\Ht) arc (180:360:0.33);
}
\foreach \i in {1,4,7} {
  \draw[aqm, draw=aqmbroken, -{Stealth[length=5.5pt,width=4.4pt]}, line width=1.0pt]
    ($ ({0.73+\i*0.66},\Ht) + (250:0.33) $) arc (250:292:0.33);
}
\draw[aqm, draw=aqmbroken, -{Stealth[length=6pt,width=5pt]}, line width=1.3pt]
  (1.7,\Ht+0.32) -- (5.5,\Ht+0.32);
\node[aqmlabsmall, anchor=south, text=aqmbroken] at (3.6,\Ht+0.40)
  {top edge: right-movers only};

% --- bottom edge: skipping orbits drifting to the left ---------------
\foreach \i in {0,1,...,9} {
  \draw[aqm, draw=aqmbroken, line width=1.0pt]
    ({0.40+\i*0.66},0) arc (180:0:0.33);
}
\foreach \i in {1,4,7} {
  \draw[aqm, draw=aqmbroken, -{Stealth[length=5.5pt,width=4.4pt]}, line width=1.0pt]
    ($ ({0.73+\i*0.66},0) + (68:0.33) $) arc (68:110:0.33);
}
\draw[aqm, draw=aqmbroken, -{Stealth[length=6pt,width=5pt]}, line width=1.3pt]
  (5.5,-0.32) -- (1.7,-0.32);
\node[aqmlabsmall, anchor=north, text=aqmbroken] at (3.6,-0.40)
  {bottom edge: left-movers only};

% --- side edges close the loop ---------------------------------------
\draw[aqm, draw=aqmbroken, -{Stealth[length=6pt,width=5pt]}, line width=1.3pt]
  (\W+0.32,\Ht-0.7) -- (\W+0.32,0.7);
\draw[aqm, draw=aqmbroken, -{Stealth[length=6pt,width=5pt]}, line width=1.3pt]
  (-0.32,0.7) -- (-0.32,\Ht-0.7);

% =====================================================================
% (b) an impurity at the edge
% =====================================================================
\begin{scope}[xshift=8.9cm, yshift=1.50cm]

\node[aqmpanel, anchor=south west] at (0,2.94) {(b) an impurity cannot stop it};

\fill[aqmfill] (0,0.10) rectangle (5.4,1.80);
\draw[aqm, draw=black] (0,1.80) -- (5.4,1.80);
\node[aqmlabsmall, anchor=south west, text=aqmmuted] at (0.06,1.90) {vacuum};
\node[aqmlabsmall, anchor=north west, text=aqmmuted] at (0.06,1.38) {sample};

% the impurity
\node[circle, draw=black, fill=aqmmuted, line width=\aqmthinwidth,
      minimum size=15pt, inner sep=0pt] at (2.70,1.68) {};
\node[aqmlabsmall, anchor=south west] at (2.98,1.98) {impurity};
\draw[aqm, draw=black, line width=\aqmthinwidth] (2.86,1.86) -- (3.02,2.02);

% the rerouted chiral mode
\draw[aqm, draw=aqmbroken, line width=1.3pt, -{Stealth[length=6pt,width=5pt]}]
  (0.85,1.52) -- (1.90,1.52)
  .. controls (2.34,1.52) and (2.26,0.95) .. (2.70,0.95)
  .. controls (3.14,0.95) and (3.06,1.52) .. (3.50,1.52)
  -- (5.10,1.52);
\node[aqmlabsmall, anchor=west, text=aqmbroken, align=left] at (3.30,0.92)
  {reroutes and\\[-1pt]carries on};

% the forbidden backscattering channel
\draw[aqm, draw=aqmmuted, line width=1.0pt, dash pattern=on 3pt off 2pt,
      -{Stealth[length=5pt,width=4pt]}] (2.90,0.45) -- (1.05,0.45);
\draw[aqm, draw=aqmbroken, line width=1.3pt] (1.75,0.24) -- (2.20,0.66);
\draw[aqm, draw=aqmbroken, line width=1.3pt] (1.75,0.66) -- (2.20,0.24);
\node[aqmlabsmall, anchor=north, align=center, text width=6.6cm] at (2.70,-0.06)
  {no counter-propagating state at this edge and this energy:
   the impurity has nothing to scatter into};

\end{scope}

\end{tikzpicture}

\end{aqmfigure}
